\exercisesection{局部上同调，Grothendieck 对偶}

\begin{enumerate}
\def\labelenumi{\arabic{enumi})}
\tightlist
\item
  设 A 为环，J 为 A 的理想，M 为 A-模。用 \(H_{A,J}(M)\) 或简作 \(H_{J}(M)\) 表示分次 A-模 \(\varinjlim_{n}\operatorname{Ext}_{A}(A/J^{n},M)\)。若 A 是局部 Noether 环，则分次 A-模 \(H_{A,m_{A}}(M)\) 与 \(H_{A}(M)\) 等同。
\end{enumerate}

\begin{enumerate}
\def\labelenumi{\alph{enumi})}
\item
  对 \(H_{J}(M)\) 建立与 \(H_{A}(M)\) 关于 A-模的同态与正合序列相类似的性质（第 1 节）。
\item
  对每个 \(i \geqslant 1\) 定义同构 \(\mathrm{H}_{\mathrm{J}}^{i+1}(\mathrm{M}) \to \varinjlim \operatorname{Ext}_{\mathrm{A}}^{i}(\mathrm{J}^{n}, \mathrm{M})\)。
\item
  假设环 A 是维数 d 的局部 Noether 环，且理想 J 由一个完全割序列 \((x_{1},\ldots,x_{r})\) 生成。构造一个同构

\[
\tau^ {i} (\mathrm{J}; \mathrm{M}): ^ {\prime} \operatorname{Tor} _ {r - i} ^ {\mathrm{A}} (\mathrm{M}, \mathrm{H} _ {\mathrm{J}} ^ {d} (\mathrm{A})) \longrightarrow \mathrm{H} _ {\mathrm{J}} ^ {i} (\mathrm{M})
\]

它推广了第 2 节的同构 \(\tau^{i}(\mathrm{M})\)；特别地，有 \(\mathrm{H}_{\mathrm{J}}^{i}(\mathrm{A}) = 0\) 对 \(i > r\) 成立。

\item
  对 A 的每个包含于 J 中的理想 K，用 \(\rho_{\mathrm{K},\mathrm{J}}:\mathrm{H}_{\mathrm{J}}(\mathrm{M})\to \mathrm{H}_{\mathrm{K}}(\mathrm{M})\) 表示由 \(\mathrm{A / K^n}\) 到 \(\mathrm{A / J^n}\) 的典范映射导出的同态。设 I 为 A 的一个理想；构造 A-模的一个长正合序列

\[
\dots \longrightarrow \mathrm{H} _ {\mathrm{I} +, \mathrm{J}} ^ {i} (\mathrm{M}) \stackrel {\alpha} {\longrightarrow} \mathrm{H} _ {\mathrm{I}} ^ {i} (\mathrm{M}) \oplus \mathrm{H} _ {\mathrm{J}} ^ {i} (\mathrm{M}) \stackrel {\beta} {\longrightarrow} \mathrm{H} _ {\mathrm{I} \cap \mathrm{J}} ^ {i} (\mathrm{M}) \longrightarrow \mathrm{H} _ {\mathrm{I} \mid \mathrm{J}} ^ {i + 1} (\mathrm{M}) \longrightarrow \dots ,
\]

\[
\text {   où   } \alpha (x) = (\rho_ {\mathrm{I,I+J}} (x), \rho_ {\mathrm{J,I+J}} (x)) \text {   et   } \beta (x, y) = \rho_ {\mathrm{I} \cap \mathrm{J}, \mathrm{I}} (x) - \rho_ {\mathrm{I} \cap \mathrm{J}, \mathrm{J}} (y).
\]

\end{enumerate}

\begin{enumerate}
\def\labelenumi{\arabic{enumi})}
\setcounter{enumi}{1}
\tightlist
\item
  设 \(\Lambda\) 为维数 \(d\) 的局部 Macaulay 环。证明为使 A 是 Gorenstein 环，必须且只需 \(\Lambda\)-模 \(\mathrm{H}_{\mathrm{A}}^{d}(\Lambda)\) 是内射的。
\end{enumerate}

¶ 3) 设 \(\rho: A \to B\) 为局部 Noether 环之间的局部同态，使得 \(\rho(\mathfrak{m}_{A})B\) 是 B 的定义理想。

\begin{enumerate}
\def\labelenumi{\alph{enumi})}
\item
  对每个 B-模 N，构造一个从 \(H_{A}(N)\) 到 \(H_{B}(N)\) 上的自然 A-线性同构。（利用 N 的一个内射分解，证明只需证明有 \(H_{A}^{i}(J)=0\) 对所有 i\textgreater0 与每个不可分解内射 B-模 J 成立。设 q 为 \(\operatorname{Ass}_{\mathrm{B}}(\mathrm{J})\) 中唯一的元素，并设 \(p=\rho^{-1}(q)\)。若 \(p\neq m_{A}\)，注意存在 \(a\in m_{A}\) 使得乘法映射 \(a_{J}\) 是双射。若 \(p=m_{A}\)，且设 I 为 A-模 J 的一个极小内射分解，注意对每个 n 有 \(\operatorname{Ass}_{\mathrm{A}}(\mathrm{I}^{n})=\{\mathfrak{m}_{\mathrm{A}}\}\)。）
\item
  进而假设 B 是平坦 A-模。对每个 A-模 M，构造一个从 \(\mathrm{B}\otimes_{\mathrm{A}}\mathrm{H}_{\mathrm{A}}(\mathrm{M})\) 到 \(\mathrm{H}_{\mathrm{B}}(\mathrm{B}\otimes_{\mathrm{A}}\mathrm{M})\) 上的自然 B-线性同构。
\end{enumerate}

\begin{enumerate}
\def\labelenumi{\arabic{enumi})}
\setcounter{enumi}{3}
\item
  设 A 为局部环，M 为非零的有限生成 A-模，维数为 e。证明 \(\mathrm{H}_{\mathrm{A}}^{e}(\mathrm{M})\) 不为零（化为环 A 完备的情形，再借助习题 3 化为 A 正则的情形）。
\item
  设 I 为 N 的有限子集。
\end{enumerate}

\begin{enumerate}
\def\labelenumi{\alph{enumi})}
\item
  构造一个正则局部环 B 与一个 B-模 N，使得满足 \(H_{\mathrm{B}}^{p}(\mathrm{N}) \neq 0\) 的整数 p 所成的集合等于 I（考虑直和）。
\item
  构造一个局部环 A，使得满足 \(\mathrm{H}_{\mathrm{A}}^{p}(\mathrm{A}) \neq 0\) 的整数 p 所成的集合等于 I（取 \(A = B \oplus N\)，其中 N 是平方为零的理想，并利用习题 3）。
\end{enumerate}

\begin{enumerate}
\def\labelenumi{\arabic{enumi})}
\setcounter{enumi}{5}
\tightlist
\item
  设 A 为具有对偶化模的局部 Noether 环，M 为有限生成 A-模，n 为整数。
\end{enumerate}

\begin{enumerate}
\def\labelenumi{\alph{enumi})}
\item
  为使 A-模 \(H_{\mathrm{A}}^{n}(\mathrm{M})\) 的长度有限，必须且只需对 A 的每个异于 \(m_{A}\) 的素理想 p，有 \(\mathrm{H}_{\mathrm{A}_{\mathfrak{p}}}^{n-\dim(\mathrm{A}/\mathfrak{p})}(\mathrm{M}_{\mathfrak{p}})=0\)（利用 Grothendieck 对偶）。
\item
  为使 \(\mathrm{H}_{\mathrm{A}}^{i}(\mathrm{M})\) 对每个 \(i \leqslant n\) 长度有限，必须且只需有 \(\operatorname{prof}(\mathrm{M}_{\mathfrak{p}}) + \dim(\mathrm{A}/\mathfrak{p}) > n\) 对每个异于 \(m_{A}\) 的素理想 p 成立
\end{enumerate}

\begin{enumerate}
\def\labelenumi{\arabic{enumi})}
\setcounter{enumi}{6}
\tightlist
\item
  设 A 为环，M 为 \(\Lambda\)-模，\(\mathbf{x} = (x_i)_{i \in I}\) 为 A 中元素的有限族。对每个整数 \(n \geqslant 1\)，用 \(\mathbf{x}^n\) 表示族 \((x_i^n)_{i \in I}\)。设 \(\Delta(\mathbf{x})\) 为 \(\mathrm{A}^{\mathrm{I}}\) 的一个自同态（由 \(\Delta(\mathbf{x})(e_i) = x_i e_i\) 定义）。
\end{enumerate}

\begin{enumerate}
\def\labelenumi{\alph{enumi})}
\tightlist
\item
  对每对整数 \((n, m)\)（满足 \(n \geqslant m\)），证明映射
\end{enumerate}

\[
\operatorname{Homgr} \left(\Lambda \left(\Delta \left(\mathbf {x} ^ {n - m}\right)\right), 1 _ {\mathrm{M}}\right): \mathbf {K} ^ {\bullet} \left(\mathbf {x} ^ {m}, \mathrm{M}\right) \longrightarrow \mathbf {K} ^ {\bullet} \left(\mathbf {x} ^ {n}, \mathrm{M}\right)
\]

定义复形的一个归纳系统；用 \(\mathbf{K}^{\bullet}(\mathbf{x}^{\infty},\mathrm{M})\) 表示这个系统的极限，用 \(\mathrm{H}^{\bullet}(\mathbf{x}^{\infty},\mathrm{M})\) 表示它的上同调。

\begin{enumerate}
\def\labelenumi{\alph{enumi})}
\setcounter{enumi}{1}
\tightlist
\item
  在 I 上选取一个全序。对 I 的每个子集 J，令 \(x_{J} = \prod_{j \in J} x_{j}\)。
\end{enumerate}

对 \(p \in \mathbf{Z}\)，令 \(\mathcal{C}^p = \bigoplus_{|J| = p} A_{x_J}\)；对 \(a \in A_{x_J}\)（其中 \(\text{Card}(J) = p\)），令 \(d^p(a) = \sum_{i \notin J} (-1)^{\varepsilon(J, i)} \frac{a}{1}\)，这里取值于 \(A_{x_{J \cup \{i\}}}\) 中，其中 \(\varepsilon(J, i)\) 是 \(J\) 中严格小于 \(i\) 的元素的个数。证明 \((\mathcal{C}, d)\) 是一个复形，并定义一个从 \(\mathbf{K}^\bullet(\mathbf{x}^\infty, M)\) 到复形 \(\mathcal{C} \otimes_A M\) 上的典范同构。

8) 沿用上一习题的记号，并进而假设环 A 是 Noether 的。用 x 表示由 x 生成的理想。

\begin{enumerate}
\def\labelenumi{\alph{enumi})}
\item
  设 m 为整数。证明存在整数 \(n_{0} \geqslant m\)，使得当 \(n \geqslant n_{0}\) 时，态射 \(\boldsymbol{\Lambda}(\Delta(\mathbf{x})^{n-m}) : \mathbf{K}_{\bullet}(\mathbf{x}^{n}, \mathrm{A}) \longrightarrow \mathbf{K}_{\bullet}(\mathbf{x}^{m}, \mathrm{A})\) 在次数 \(\geqslant 1\) 的同调上诱导 0（注意有 \(\boldsymbol{\Lambda}^{p}(\Delta(\mathbf{x}^{n-m}))(\mathbf{K}_{p}(\mathbf{x}^{n}, \mathrm{A})) \subset x^{n-m}\mathbf{K}_{p}(\mathbf{x}^{m}, \mathrm{A})\)；由 Artin–Rees 引理推出存在整数 \(n_{0}\)，使 \(Z_{p}(\mathbf{x}^{m}, \mathrm{A}) \cap x^{n_{0}}\mathbf{K}_{\bullet}(\mathbf{x}^{m}, \mathrm{A})\) 包含于 \(x^{rm}Z_{p}(\mathbf{x}^{m}, \mathrm{A})\)，从而包含于 \(B_{p}(\mathbf{x}^{m}, \mathrm{A})\)）。
\item
  定义一个从 \(\varinjlim_{n}\mathrm{Ext}_{\mathrm{A}}^{p}(\mathrm{A}/x^{n},\mathrm{M})\) 到 \(\mathrm{H}^p (\mathbf{x}^\infty ,\mathrm{M})\) 上的典范同构（化为证明 \(\mathrm{H}^p (\mathbf{x}^\infty ,\mathrm{M})\) 在 \(p\geqslant 1\) 且 M 内射时为零；在这一情形注意 \(\mathrm{H}^p (\mathbf{x}^\infty ,\mathrm{M})\) 与 \(\mathrm{Hom}_{\mathrm{A}}(\mathrm{H}_p(\mathbf{x}^n,\mathrm{M}))\) 的归纳极限等同，并利用 \(a)\)）。特别地，若 A 是局部的且 \(x\) 是 A 的定义理想，则分次 A-模 \(\mathrm{H}_{\mathrm{A}}(\mathrm{M})\) 典范同构于 \(\mathrm{H}^{\bullet}(\mathbf{x}^{\infty},\mathrm{M})\)。
\item
  假设环 A 是维数 \(d\) 的局部 Noether 环；设 \(\mathbf{x} = (x_{1},\ldots ,x_{d})\) 为 \(\mathfrak{m}_{\Lambda}\) 中元素的一个极大割序列。证明 A-模 \(\mathrm{H}_{\mathrm{A}}^{d}(\mathrm{M})\) 与 A-模的归纳系统 \((\mathrm{M}_p,u_{qp})\) 的极限等同，其中 \(\mathrm{M}_p\) 为商模 \(\mathrm{M} / (x_1^p\mathrm{M} + \dots +x_d^p\mathrm{M})\)，而 \(u_{qp}:\mathrm{M}_p\to \mathrm{M}_q\)（\(q\geqslant p\)）是由比率为 \((x_{1}\dots x_{d})^{q - p}\) 的乘法映射诱导的同态。
\end{enumerate}

\begin{enumerate}
\def\labelenumi{\arabic{enumi})}
\setcounter{enumi}{8}
\tightlist
\item
  设 A 为 Noether 局部环，d 为其维数，\(\mathbf{x} = (x_{1}, \ldots, x_{d})\) 为 \(m_{A}\) 中元素的一个极大割序列。
\end{enumerate}

\begin{enumerate}
\def\labelenumi{\alph{enumi})}
\item
  证明存在整数 \(m_0\)，使得当 \(m \geqslant m_0\) 且 \(n \geqslant 0\) 时，有 \((x_1^m \ldots x_d^m)^n \not\in (x_1^{m(n+1)}, \ldots, x_d^{m(n+1)})\)（利用习题 4 与 7, b)）。
\item
  假设 \(\Lambda\) 包含一个特征为 \(p\) 的域。证明有 \((x_{1} \ldots x_{d})^{n} \notin (x_{1}^{n+1}, \ldots, x_{d}^{n+1})\) 对每个 \(n \geqslant 0\) 成立，即 A 满足条件 (CM)（§ 2 习题 3；考察自同态 \(x \mapsto x^p\)）。因此，每个包含域的 Noether 局部环都满足条件 (CM) \(^{1}\)。
\item
  再假设 A 是正则的；设 \(\rho: A \to B\) 为一个单的环同态，使 B 成为有限生成的 A-模。证明 A-模 \(\rho(A)\) 是 B 的直因子（参见 § 2, 习题 3）；化为 B 是整环从而是局部的情形，并应用 § 4 习题 8）。
\end{enumerate}

\begin{enumerate}
\def\labelenumi{\arabic{enumi})}
\setcounter{enumi}{9}
\tightlist
\item
  设 \(k\) 为域，A 为 \(k\) 上的维数 \(d\) 的正则局部代数，且扩张 \(k \to \kappa_{\mathrm{A}}\) 是平凡的。A-模 \(\mathrm{H}_{\mathrm{A}}^{d}(\mathrm{A})\) 与 \(\mathrm{A}' = \mathrm{Hom}_k^{cont}(\mathrm{A}, k)\) 都是 Matlis 模，从而同构；我们要在这两个模之间把一个同构明显地写出来。
\end{enumerate}

\begin{enumerate}
\def\labelenumi{\alph{enumi})}
\item
  设 \(\mathbf{x} = (x_{1}, \ldots, x_{d})\) 为 A 的一个坐标系。对每个整数 \(p \geqslant 0\)，用 \(A_{p}\) 表示 k-代数 \(\mathrm{A}/(x_{1}^{p}, \ldots, x_{d}^{p})\)。单项式 \(x_{1}^{a_{1}} \ldots x_{d}^{a_{d}}\)（其中 \(0 \leqslant a_{i} < p\)）的类构成 \(A_{p}\) 在 k 上的一组基；用 \(\rho_{p}\) 表示 \(A_{p}\) 上的线性形式，它在 \(x_{1}^{p-1} \ldots x_{d}^{p-1}\) 的类上取值 1，在基的其余元素上取值 0。证明 k-线性映射 \(\tilde{\rho}_{p}: A_{p} \to \operatorname{Hom}_{k}(A_{p}, k)\)（它由双线性形式 \((a, b) \mapsto \rho_{p}(ab)\) 导出）是同构（可直接论证，或利用 § 5 习题 10）。
\item
  对 \(q \geqslant p\)，设 \(u_{qp}: A_p \to A_q\) 为由比率为 \((x_1 \ldots x_d)^{q-p}\) 的乘法映射诱导的 A-线性同态，并设 \(\pi_{pq}: A_q \to A_p\) 为商映射；证明有 \(\tilde{\rho}_q \circ u_{qp} = {}^t \pi_{pq} \circ \tilde{\rho}_p\)。取归纳极限并利用习题 8, c)，由此得到一个 A-线性同构 \(\tilde{\rho}: H_A^d(A) \to A'\)。
\item
  假设 \(d=1\)，于是 A 是离散赋值环，且 \(x_{1}=x\) 是 A 的一个单值化元。证明单值化元 \(x\) 的选取使 \(\mathrm{H}_{\mathrm{A}}^{1}(\mathrm{A})\) 得以与 K/A 等同；同构 \(\rho:\mathrm{K}/\mathrm{A}\to\mathrm{A}'\) 于是由 \(\rho(\varphi)(a)=\mathrm{Res}(a\varphi)\) 定义，其中 \(\mathrm{Res}:\mathrm{K}/\mathrm{A}\to k\) 是到因子 \(k x^{-1}\)（在典范分解 \(\mathrm{K}/\Lambda=\bigoplus_{n\geqslant 1}k x^{-n}\) 中）上的投射。
\end{enumerate}

\begin{enumerate}
\def\labelenumi{\arabic{enumi})}
\setcounter{enumi}{10}
\item
  设 A 为 Noether 局部环，I 为 A 上的一个 Matlis 模，M 为有限生成 A-模。令 \(\mu^{p} = \dim_{\kappa_{A}} \operatorname{Ext}_{A}^{p}(\kappa_{A}, M)\) 对每个 \(p \geqslant 0\)（参见 § 8, 习题 4）。证明分次 A-模 H \(_{A}\) (M) 与一个复形 \(\ldots \to 0 \to I^{\mu^{0}} \to I^{\mu^{1}} \to \ldots\) 的同调等同（见该处）。特别地，若 M 非零，则有 \(\mu^{p} \neq 0\) 对 \(p = \text{prof}(M)\) 与 \(p = \dim(M)\)（习题 4）。
\item
  设 A 为 Noether 局部环，M 为维数 d 的有限生成 A-模；假设 M 是 Buchsbaum 模（§ 2, 习题 7）。
\end{enumerate}

\begin{enumerate}
\def\labelenumi{\alph{enumi})}
\item
  证明 \(\mathrm{H}_{\mathrm{A}}^{i}(\mathrm{M})\) 被 \(m_{A}\) 零化（对 \(i \neq d\)）（先处理 d = 1 的情形，再对 d 作归纳）。
\item
  设 \(x\) 为关于 \(\mathfrak{m}\) 的一个割元素；对 \(0 \leqslant i \leqslant d - 2\) 定义正合序列 \(0 \to \mathrm{H}_{\mathrm{A}}^{i}(\mathrm{M}) \to \mathrm{H}_{\mathrm{A}}^{i}(\mathrm{M}/x\mathrm{M}) \to \mathrm{H}_{\mathrm{A}}^{i+1}(\mathrm{M}) \to 0\)。
\item
  证明公式 \(i(\mathrm{M}) = \sum_{i=0}^{d-1} \left( \begin{array}{cc} d & 1 \\ i & \end{array} \right) \lg(\mathrm{H}_{\mathrm{A}}^i(\mathrm{M}))\)（对整数 \(d \geqslant 1\) 作归纳，利用 \(b\) ) 与 § 2 习题 8, \(d\)）。
\end{enumerate}

